\subsection{借紧算子的扰动}

定理 2. --- 设 E 是完备可赋范空间，u 是 E 的一个自同态，h 是 E 的一个紧自同态。则 \(\mathrm{Sp}_{\omega}(u + h) = \mathrm{Sp}_{\omega}(u)\) ，且 \(\mathrm{Sp}_n(u + h) = \mathrm{Sp}_n(u)\) （对一切 \(n \in \overline{\mathbf{Z}} - \{0\}\) ）。

设 \(\lambda \in \mathbf{C}\) 。为使 \(u + h - \lambda 1_{\mathrm{E}}\) 是核（相应地，余核）有限维的严格态射，必须且只需 \(u - \lambda 1_{\mathrm{E}}\) 具有同样的性质，这由 III, p. 72 的定理 1（相应地，III, p. 73 的定理 2）得出。等式

\[
\begin{array}{c} \operatorname{Sp} _ {- \infty} (u + h) = \operatorname{Sp} _ {- \infty} (u), \quad \operatorname{Sp} _ {+ \infty} (u + h) = \operatorname{Sp} _ {+ \infty} (u), \\ \operatorname{Sp} _ {\omega} (u + h) = \operatorname{Sp} _ {\omega} (u) \end{array}
\]

由此得出。此外，我们有 \(\operatorname{Sp}_{n}(u+h)=\operatorname{Sp}_{n}(u)\) （由 III, p. 73 的定理 3，对每个 \(n\in \mathbf{Z}-\{0\}\) ）。

推论. --- 假设 \(\mathrm{Sp}_{\omega}(u)\) 的补集的无界连通分支包含 0。则 \(u + h\) 是 E 的里斯自同态。

有 \(\mathrm{Sp}_{\omega}(u+h)=\mathrm{Sp}_{\omega}(u)\) （由定理 2），故 0 属于 \(\mathbf{C}-\mathrm{Sp}_{\omega}(u+h)\) 的无界连通分支。由 III, p. 87 的定理 1 的推论，或者 0 不属于 \(u+h\) 的谱，或者它是该谱的敏感点。在两种情形下，\(u+h\) 都是 E 的里斯自同态。

